\subsection{可解圈空间}

定义 2. --- 称拓扑空间 B 为可解圈的，如果它是局部道路连通的，并且 B 的每一点 b 都有一个邻域 V，使得由典范单射导出的、从 \(\pi_1(V,b)\) 到 \(\pi_1(B,b)\) 中的同态以单位元为其像。

一个局部道路连通的拓扑空间是可解圈的，当且仅当它的每个连通分支都是可解圈的。

设 B 是一个局部道路连通的拓扑空间。假设 B 的每一点都有一个可解圈的邻域 V。那么 B 是可解圈的。

每个局部道路连通且道路单连通的拓扑空间都是可解圈的。同伦于一点的局部道路连通空间尤其如此。

注. --- 1) 存在连通且局部道路连通的拓扑空间 B，使得 B 的某些点 a（但并非所有点）具有邻域 V，使从 \(\pi_{1}(V,a)\) 到 \(\pi_{1}(B,a)\) 中的同态是平凡的（III, p. 336, 习题 6）。

\begin{enumerate}
\def\labelenumi{\arabic{enumi})}
\setcounter{enumi}{1}
\tightlist
\item
  设 B 是一个可解圈空间；群胚 \(\varpi(\mathrm{B})\) 在 B-空间上的每个作用都无局部单值（参见 III, p. 313, 注）。特别地，为使映射 \(p: E \rightarrow B\) 使 E 成为 B 的覆叠，必要且充分的是：它平展、分离，并且满足道路提升性质（III, p. 315, 推论 3）。
\end{enumerate}

命题 1. --- 有限多个可解圈空间的乘积空间是可解圈的。

只需证明：若 A 与 B 是两个可解圈空间，则其乘积 A × B 也是。在这些条件下，空间 A × B 确实是局部道路连通的（III, p. 261, 命题 9）。设 (a, b) 是 A × B 的一点。依假设，存在 a 的一个邻域 U（相应地，b 的一个邻域 V），使同态 \(i_* \colon \pi_1(\mathrm{U}, a) \to \pi_1(\mathrm{A}, a)\)（它由典范单射 \(i \colon \mathrm{U} \to \mathrm{A}\) 导出；相应地，同态 \(j_* \colon \pi_1(\mathrm{V}, b) \to \pi_1(\mathrm{B}, b)\) 由典范单射 \(j \colon \mathrm{V} \to \mathrm{B}\) 导出）的像退缩为单位元。于是 U × V 是 (a, b) 在 A × B 中的一个邻域。同态 \(\pi_1(\mathrm{U} \times \mathrm{V}, (a, b)) \to \pi_1(\mathrm{A} \times \mathrm{B}, (a, b))\) 与同态 \((i_*, j_*)\) 等同（III, p. 297, 命题 4 的推论）。因此其像退缩为单位元。这就证明了 A × B 是可解圈的。

命题 2. --- a) 可解圈空间的每个覆叠都是可解圈的。

\begin{enumerate}
\def\labelenumi{\alph{enumi})}
\setcounter{enumi}{1}
\item
  反之，设 \(p \colon \mathrm{E} \to \mathrm{B}\) 是一个平展且满射的映射，并设 \(\mathrm{E}\) 是可解圈空间。那么 \(\mathrm{B}\) 是可解圈的。
\item
  设 B 是一个可解圈拓扑空间，\((\mathrm{E}, p)\) 是 B 的一个覆叠。于是空间 E 局部道路连通。设 x 是 E 的一点，记 \(b = p(x)\)，并设 V 是 b 在 B 中的一个邻域，使典范同态 \(\pi_{1}(\mathrm{V}, b) \to \pi_{1}(\mathrm{B}, b)\) 平凡。用 \(i: V \to B\) 与 \(j: p^{-1}(V) \to E\) 表示典范单射。依假设，同态 \(\pi_{1}(i, b)\) 的像退缩为单位元。由于 \(p \circ i = j \circ p\) 且 \(\pi_{1}(p, x)\) 是单射的，同态 \(\pi_{1}(j, x)\) 的像退缩为单位元。这就证明了 E 是可解圈的。
\item
  空间 B 局部道路连通。设 b 是 B 的一点，x 是 \(E_{b}\) 的一点。由于 p 平展，存在 x 在 E 中的一个邻域 U，使 p 诱导 U 到 \(p(\mathrm{U})\) 上的一个同胚。设 V 是 x 在 E 中的一个包含于 U 的邻域，使同态 \(\pi_{1}(j,x)\) 平凡，其中 \(j:V\to E\) 表示典范单射。用 \(q:V\to p(V)\) 表示由 p 取子空间而导出的映射；它是一个同胚。再用 i 表示 \(p(\mathrm{V})\) 到 B 中的典范单射。我们有 \(p\circ j=i\circ q\)，因此 \(\pi_{1}(i,b)\circ\pi_{1}(q,x)=\pi_{1}(p,x)\circ\pi_{1}(j,x)\) 是平凡同态。由于同态 \(\pi_{1}(q,x)\) 是同构，同态 \(\pi_{1}(i,b)\) 平凡。由此得出空间 B 是可解圈的。
\end{enumerate}

命题 3. --- 设 B 是一个可解圈拓扑空间。设 (Y, q) 是 B 的一个覆叠，(Z, p) 是 Y 的一个覆叠。于是赋以映射 \(q \circ p\) 的拓扑空间 Z 是 B 的一个覆叠。

事实上，映射 \(q \circ p\) 是平展的（I, p. 29, 命题 6）且分离的（I, p. 25, 命题 2）。由上面的注 2，因此只需证明它满足道路提升性质。设 \(z\) 是 Z 的一点，并设 \(c \colon \mathbf{I} \to \mathbf{B}\) 是 B 中以 \(q(p(z))\) 为起点的一条道路。存在一条道路 \(c'\)，它以 \(p(z)\) 为起点、位于 Y 中，且提升 \(c\)，因为 Y 是 B 的覆叠（III, p. 302, 推论 2）。由于 Z 是 Y 的覆叠，存在一条道路 \(c''\)，它以 \(z\) 为起点、位于 Z 中，且提升 \(c'\)；道路 \(c''\) 提升 \(c\)。命题得证。

